% Research continuity report. Self-contained: no external bibliography or figures.
% Build: pdflatex batched_mpc_oram_report.tex (twice), or latexmk -pdf.
% Restart points: Sections 1, 12, and 13; estimates are conditional, not benchmarks.
\documentclass[11pt]{article}
\usepackage[letterpaper,margin=0.88in]{geometry}
\usepackage[T1]{fontenc}
\usepackage{lmodern,microtype}
\usepackage{amsmath,amssymb,mathtools}
\usepackage{booktabs,tabularx,longtable,array}
\usepackage{xcolor,enumitem,fancyhdr}
\usepackage{listings}
\usepackage[hyphens]{url}
\usepackage{hyperref}
\definecolor{navy}{HTML}{17324D}
\definecolor{teal}{HTML}{176B73}
\hypersetup{colorlinks=true,linkcolor=navy,citecolor=teal,urlcolor=teal,
 pdftitle={Batched Oblivious Lookup in Ten-Party MPC},
 pdfauthor={Research discussion consolidated with ChatGPT}}
\setlength{\parindent}{0pt}
\setlength{\parskip}{5pt plus 1pt}
\setlength{\emergencystretch}{2em}
\setlist{nosep,leftmargin=1.5em,topsep=5pt}
\renewcommand{\arraystretch}{1.18}
\newcolumntype{Y}{>{\raggedright\arraybackslash}X}
\newcommand{\F}{\mathbb F}
\newcommand{\sh}[1]{\left[#1\right]_4}
\newcommand{\E}{\mathbb E}
\newcommand{\status}[1]{\textbf{Status: #1.}}
\newcommand{\term}[1]{\textbf{#1}}
\newcommand{\smallurl}[1]{{\small\url{#1}}}
\pagestyle{fancy}
\fancyhf{}
\fancyhead[L]{\small\textcolor{navy}{Batched oblivious lookup in MPC}}
\fancyhead[R]{\small Research continuity report}
\fancyfoot[C]{\thepage}
\setlength{\headheight}{14pt}
\lstset{basicstyle=\ttfamily\small,breaklines=true,columns=fullflexible,
 frame=single,rulecolor=\color{black!20},xleftmargin=3pt,xrightmargin=3pt}

\begin{document}
\hypersetup{pageanchor=false}
\begin{titlepage}
\thispagestyle{empty}
\vspace*{0.6in}
{\color{teal}\large RESEARCH CONTINUITY REPORT\par}
\vspace{0.22in}
{\color{navy}\Huge\bfseries Batched Oblivious Lookup\\[7pt] in Ten-Party MPC\par}
\vspace{0.28in}
{\Large Square-root and tree ORAM, batched array access,\\[4pt]
honest-majority preprocessing, equality, and deployment costs\par}
\vspace{0.38in}
{\large Consolidated from the research discussion\\13 October 2026 (UTC)\par}
\vspace{0.42in}
\rule{\textwidth}{0.6pt}
\begin{abstract}
\noindent This report preserves the problem statement, candidate constructions,
literature, derivations, corrections, and unresolved questions from a discussion
of a database containing $2^{28}$ pairs of 32-byte keys and 32-byte values.
The target is 10,000 read requests per second across ten MPC parties, using
batches of approximately 1,024 requests and parallel workers. The security
assumption evolved from conventional SPDZ-style dishonest-majority computation
to allowing an honest majority. Degree-4 Shamir sharing among ten parties is
therefore the main arithmetic representation considered here.

The central research direction is to combine efficient batched private array
access with an enlarged square-root ORAM shelter and a shallow recursive
position map. The report derives when this could outperform a direct batched
scan, explains why compaction alone is insufficient, and records tree-based
alternatives. It also gives explicit equality protocols and a conditional AWS
cost model. No complete ten-party ORAM implementation or end-to-end throughput
measurement was produced in the discussion. Numerical projections must retain
their stated assumptions.
\end{abstract}
\vfill
\textbf{Reading guide.} Start with Section~\ref{sec:restart}. The algorithmic
models are in Sections~\ref{sec:batch}--\ref{sec:trees}; equality is in
Section~\ref{sec:equality}; cost accounting is in Section~\ref{sec:cost}.
Section~\ref{sec:next} identifies the next experiments.
\end{titlepage}

\pagenumbering{roman}
\hypersetup{pageanchor=true}
\tableofcontents
\clearpage
\pagenumbering{arabic}

\section{Restart summary and status of conclusions}\label{sec:restart}

\textbf{Working direction.} Prototype an honest-majority arithmetic backend
with batched gather/merge primitives before selecting the ORAM. Compare a direct
batched array access, a recursive square-root ORAM with an enlarged shelter,
and a tree ORAM with parallelized metadata processing using the same field,
security threshold, payload encoding, and complete preprocessing costs.

The following points should survive into later sessions.
\begin{enumerate}
\item \textbf{The user accepts honest majority.} Ten parties with degree-4
Shamir sharing suggests privacy against up to four corruptions. Malicious
security with abort is the working interpretation, not an explicitly finalized
deployment requirement. Do not silently substitute a three-corruption protocol
requiring $t<n/3$, or assume passive security is sufficient.
\item \textbf{The square-root proposal is promising but unbenchmarked.}
An idealized shelter/rebuild model gives $S\asymp\sqrt{Db}$, where $D$ is dense
array capacity and $b$ is the number of dense accesses in a batch. For one dense
access per request, $D=2^{28}$ and $b=2^{10}$ give $S=2^{19}$.
This is a derived balance, not a measured optimum.
\item \textbf{The fair baseline is a batched gather.} It need not perform
$Db$ equalities. Private stable sorting/routing, finite differences, and local
prefix sums give a substantially better baseline. Compaction requires selection
marks; it does not create membership predicates for free.
\item \textbf{Duplicate handling depends on the construction.} A direct
batched read can support duplicates without an explicit deduplication pass.
An ORAM batch must still avoid exposing repeated physical probes or hit counts.
Recursive position-map accesses can duplicate even when data addresses differ.
\item \textbf{Equality need not decompose both inputs.} A fresh random field
mask with shared canonical bits yields an exact secret equality bit. A simple
product-tree implementation takes $\ell-1$ secure multiplications after one
opening for an $\ell$-bit prime field. A separately derived power-tuple variant
moves this product work offline and uses two dependent online openings.
\item \textbf{Honest majority changes the preprocessing economics.}
Conventional dishonest-majority LowGear/SPDZ triple rates are not the correct
baseline. ATLAS and Shamir-based malicious protocols are more relevant, but
published synthetic-circuit rates are not ORAM throughput measurements.
\item \textbf{No final annual budget has been established.} The earlier
rough \$0.3--1 billion/year dishonest-majority projection is superseded as a
baseline. The honest-majority model remains conditional on the as-yet-unmeasured
work per lookup and deployment topology. Cross-AZ bandwidth can itself dominate.
\end{enumerate}

\subsection{Evidence labels}

\begin{tabularx}{\textwidth}{@{}lY@{}}
\toprule
Label & Meaning in this report \\
\midrule
Requirement & Explicitly supplied by the user. \\
Working assumption & Chosen to make an analysis concrete; still subject to design decisions. \\
Published result & Attributed to an identified paper or maintained implementation. \\
Derived model & Algebra or a construction assembled in the discussion, with assumptions stated. \\
Unresolved & Requires a proof, implementation, parameter choice, or measurement. \\
\bottomrule
\end{tabularx}

The report is a technical continuity record, not a security proof for a complete
system. Absence of a located paper with the exact proposed combination is not
evidence of novelty. Literature coverage reflects the discussion and selected
checks through the consolidation date.

\section{Workload, notation, and security model}\label{sec:workload}

\begin{tabularx}{\textwidth}{@{}l l Y@{}}
\toprule
Symbol & Value or role & Interpretation \\
\midrule
$N$ & $2^{28}$ & Number of logical key--value pairs. \\
Key/value & $32+32$ bytes & Sparse 256-bit keys, 256-bit values. \\
$Q$ & $1024=2^{10}$ & Logical map requests per batch. \\
$R$ & $10{,}000\ \mathrm{s}^{-1}$ & Target aggregate logical read rate. \\
$n,t$ & $10,4$ & Parties and working corruption threshold. \\
$D$ & Construction-dependent & Number of dense cells after map representation. \\
$h$ & Usually 1 or 2 & Dense accesses per logical map request. \\
$b$ & $hQ$ & Dense accesses per batch. \\
$S$ & Tunable & Shelter capacity / padded accesses per epoch. \\
$F_j$ & Tunable & Number of position labels packed into a recursive block. \\
$w$ & Context-dependent & Field-word width in bytes, or payload width when explicitly stated. \\
$G$ & Unmeasured & Full-backend field multiplications per logical lookup. \\
\bottomrule
\end{tabularx}

The raw database size is
\[
N(32+32)=2^{34}\ \text{bytes}=16\ \text{GiB}.
\]
This excludes hash-table slack, ORAM buckets, dummy cells, position maps,
authentication, field-encoding overhead, shelter storage, and worker copies.
For example, a dense table with $D\simeq 2N$ and 64-byte cells already occupies
32 GiB before those additions. Field shares are not generally packed as the
original bytes: encoding a 256-bit value in five 61-bit field elements consumes
40 bytes per party before any tags.

\subsection{Batch semantics and latency}

The main analysis assumes \emph{nonadaptive batches}: all addresses in a batch
are available before its evaluation. Dependent traversals within one lookup
still require sequential phases. Batching independent lookups fills circuit
layers and amortizes overhead; it does not remove a data dependency.

At 10,000 requests/s, a single stream takes $Q/R=102.4$ ms to accumulate 1,024
requests. Uniform arrivals give roughly 51.2 ms average batch-formation waiting
time before computation. Splitting traffic across many independent workers
increases this delay unless batches are formed centrally, smaller batches are
used, or sufficient additional concurrency is available. This matters when
combining a throughput target with a future latency objective.

Application writes are out of scope for the cost target. \emph{ORAM state still
changes on reads}: locations are remapped, shelters updated, and arrays shuffled
or rebuilt. Immutable logical data makes replication easier but does not make
an individual ORAM replica stateless.

\subsection{What honest majority does and does not fix}

Degree-4 Shamir sharing among ten parties permits the usual honest-majority
multiplication regime because $2t=8<n=10$. This is a statement about sharing and
the availability of suitable protocols. Active correctness requires verified
preprocessing, consistency checks, and an appropriate opening procedure.

The desired leakage should include public database capacity, batch size, and a
fixed service schedule, but exclude logical addresses, keys, equality outcomes,
hit/miss status, duplicate structure, and the number of distinct requests unless
explicitly allowed. Aborts and output release must follow the chosen malicious
protocol. Using one AWS provider or one availability zone is an operational
assumption; it is not a proof of independence between the parties.

\section{A common accounting framework}\label{sec:accounting}

For each candidate, report at least the vector
\[
(M_{\rm field},M_{\rm bit},C_{\rm eq},C_{\rm order},
 W_{\rm perm},B_{\rm comm},D_{\rm interactive},M_{\rm RAM}),
\]
where the entries denote arithmetic work, Boolean work, equality and ordering
operations, words privately permuted, bytes communicated, interactive depth,
and memory. A scalar cost is useful only after assigning backend-specific
weights. In particular, ordering a secret address and moving one payload word
are different operations.

Conditional movement of one field word has the form
\[
\operatorname{mux}(s,x,y)=y+s(x-y),
\]
so a selector applied to a multiword record generally entails multiple secret
products. Reusing the selector may enable vectorized authentication or checks,
but the record width still matters. Conversely, moving local share arrays by a
public index costs no MPC multiplication. The cost of a \emph{private}
permutation depends on how its control information is represented and prepared.

Maintain four separate ledgers:
\begin{enumerate}
\item online work and dependent communication phases;
\item preprocessing production, verification, and storage;
\item amortized rebuild and position-map maintenance;
\item steady-state wall time, RAM traffic, network traffic, and cloud charges.
\end{enumerate}
An online-only speedup is insufficient for a continuously operating service if
preprocessing cannot be replenished at the same rate. An annual estimate must
also account for the cost of all ten parties rather than just one.

\section{Direct batched access, compaction, and duplicate requests}\label{sec:batch}

\subsection{The direct baseline is not a batch of independent scans}

Performing $b$ independent scans of $D$ records costs $\Theta(Db)$ predicates.
That is a useful simple implementation, but an unfair sole baseline for the
proposed batched ORAM. Laud's arithmetic-black-box array access gives a stronger
alternative using private stable sorting and permutation operations
\cite{laud2015}. Its read construction can be expressed as follows.

Prepare a stable sorting permutation $\pi$ for
\[
(0,1,\ldots,D-1,a_0,\ldots,a_{b-1}),
\]
where database indices are public and query indices are secret. Form
\[
\Delta X=(X_0,X_1-X_0,\ldots,X_{D-1}-X_{D-2}),
\]
append $b$ zero records, apply $\pi$, take prefix sums locally on the shares,
and apply $\pi^{-1}$. The query positions contain their requested values because
\[
\sum_{i\leq a}\Delta X_i=X_a.
\]
Stability puts a database row before its equal-address queries. Repeated queries
add zeros, so all obtain the same answer. This read method does not require a
separate query deduplication step or secret conditional copying of every
payload word in a segmented scan.

For accounting, write
\begin{equation}
F(D,b,w)=\operatorname{PrepareRead}(D,b)
+2\operatorname{ApplyPerm}(D+b,w)+O(Dw)_{\rm local}.
\label{eq:direct}
\end{equation}
The preparation term is essential. A generic oblivious sorting network may
dominate it. More specialized stable merge/radix/shuffle constructions may be
better, depending on the MPC backend. Computing a narrow permutation before
applying it to wide records avoids charging the entire payload width at every
metadata-sorting stage. The map from 256-bit keys to dense indices is a separate
problem; this baseline is immediately applicable to an indexed array and to
recursive position maps.

\subsection{What compaction actually contributes}

Oblivious compaction accepts secret marks and moves marked records to designated
positions without revealing the marks. It does not decide which records match
the secret requests. If producing each mark still compares a record with every
query, the $Db$ work remains. A meaningful claimed improvement must provide the
complete private join, routing, or membership procedure.

Stable compaction, merging, and secure shuffles are therefore relevant building
blocks \cite{falk2021,blunk2025,logstar,shuffle2024}, but they have different
party counts, security assumptions, online/offline splits, and output formats.
A two-party or passive three-party timing cannot be inserted into the ten-party
malicious model without an implementation and security adaptation.

Stable three-party sorting \cite{sorting2022} and sorting-based private
JOIN/GROUP-BY \cite{joins2024} provide relevant database precedents for separating
narrow control from wide records. Linear-work tight compaction is also possible
in an oblivious parallel model \cite{compaction2020}; neither an inherent
$N\log N$ compaction cost nor free matching should be assumed. Stability,
cardinality leakage, and output multiplicity depend on the chosen primitive.

\subsection{A sorted-shelter extension}\label{sec:sorted-shelter}

\status{Derived construction idea}
Suppose the shelter contains valid records sorted by secret logical address.
Stable-merge its addresses with the batch addresses, placing a shelter record
before equal queries. Apply the same finite-difference and prefix-sum idea to
the shelter's address, value, and validity columns. After inverse routing, each
query receives its predecessor's columns. Its membership bit is
\[
\operatorname{hit}_j=
\operatorname{predValid}_j\cdot
\operatorname{EQ}(\operatorname{predAddr}_j,a_j).
\]
This needs only $b$ final equalities, rather than $Sb$ equality tests. A validity-0
sentinel below the valid address domain handles queries with no predecessor;
invalid padding must sort after valid addresses. Sentinel and padding order
should use an explicit tagged representation rather than an accidental field
ordering convention.

The unresolved cost is the private stable merge permutation, its application
and inverse, and maintenance of the sorted shelter. A comparator-based merge
can still contain many expensive secret ordering operations. Thus the extension
does not establish linear cost with a small constant on the target backend.

\subsection{Deduplication and inverse routing}

An ORAM needs more than a functional batched read: the publicly observed physical
accesses must not expose repeats or shelter hits. A typical batch wrapper sorts
the query addresses, forms secret group-boundary bits, chooses representatives,
performs the real operations with fixed-size cover traffic, copies answers to
duplicates, and restores original order. Counts and group structure remain
secret. All of these steps belong in the cost.

For scale, a bitonic sorter on $b=1024$ inputs has
\[
\frac{b}{4}\log_2 b(\log_2 b+1)=28{,}160
\]
compare-exchanges and 55 comparator layers. This is a metadata reference count,
not a recommendation to use bitonic sorting, and not an arithmetic multiplication
count. For $b=2048$, the corresponding counts are 67,584 and 66.

With read-only data and a unique shelter, inserting only distinct real misses
avoids a full-shelter adjacent-equality pass: these addresses are absent from
the old shelter and mutually distinct. There is a qualification if cover probes
fetch unused \emph{real} rows instead of explicit dummy blocks. Such rows must
also be cached in the shelter, or made available through another equally
oblivious mechanism, so future requests need not revisit their exposed physical
locations. Fresh real cover probes must also avoid the batch's other fresh
probes, including true misses. Cover probes, duplicate suppression, and
dummy capacity need one coherent security argument.

\section{Square-root ORAM: balance and comparison with direct access}\label{sec:sqrt}

\subsection{Construction skeleton and known starting point}

The proposal uses a shuffled main array, a shelter of touched records, and a
fixed rebuild schedule. A query is resolved privately against the shelter;
otherwise it accesses the appropriate as-yet-unexposed main-array location.
Cover accesses hide which case occurred. The position map is itself recursively
protected, with a flat batched-access base case. Zahur et al.'s
\emph{Revisiting Square-Root ORAM} is a concrete MPC-oriented starting point
with recursive permutation maps and bulk initialization \cite{zahur2016}.

The particular combination proposed here---enlarged shelter, batched private
shelter merge, high position-map fanout, and a large base case on an
honest-majority backend---has not been implemented or fully proved in this
discussion. The following analysis is conditional on realizing its component
costs and access schedule securely.

\subsection{General amortization model}

Let $J(S,b)$ be the \emph{complete} batch shelter operation, including predicates,
routing, duplicate handling, output propagation, and maintenance as appropriate.
Let $R_{\rm build}(D)$ be a complete main-array rebuild including its position-map
work. After a public schedule of approximately $S/b$ batches, rebuilding yields
\begin{align}
C_{\rm sqrt,batch}(S)&=J(S,b)+\frac{b}{S}R_{\rm build}(D),\label{eq:sqrt-general}\\
C_{\rm sqrt,map}(S)&=\frac{J(S,b)}{Q}+\frac{h}{S}R_{\rm build}(D).
\end{align}
Extra fixed per-access work and recursion must be added. The epoch must be
scheduled by public capacity, not by revealing the number of distinct misses.
These equations use one consistent scalar resource; they cannot mix online
latency in $J$ with total communication in $R_{\rm build}$.

Suppose the measured range admits
\[
J(S,b)\simeq c_HS+H(b),\qquad R_{\rm build}(D)\simeq c_RD.
\]
Then differentiating gives
\begin{align}
S^*&=\sqrt{\frac{c_R}{c_H}Db},\label{eq:s-opt}\\
C_{\rm sqrt,batch}^*&=2\sqrt{c_Hc_RDb}+H(b).\label{eq:sqrt-opt}
\end{align}
Thus $S\asymp\sqrt{Db}$ is the equal-coefficient specialization. If rebuilding
costs $\Theta(D\log D)$ while shelter work is linear, the balance instead becomes
$\Theta(\sqrt{Db\log D})$. With logarithmic shelter routing, optimize the actual
expression; in general the stationary condition is
\[
\frac{\partial J(S,b)}{\partial S}
=\frac{bR_{\rm build}(D)}{S^2}.
\]
For whole-batch epochs, require $S\geq b$ and a public batch-aligned schedule;
round or project the stationary optimum into the construction's feasible range.
Include material dummy padding in the actual rebuilt capacity. If the stationary
point is infeasible, its formula is only an unconstrained lower value for the
model, and the constrained cost must be evaluated instead.

\begin{table}[htbp]
\centering
\begin{tabular}{@{}lrrrr@{}}
\toprule
Illustrative layout & $D$ & $b$ & $S=\sqrt{Db}$ & Raw shelter \\
\midrule
One dense access & $2^{28}$ & $2^{10}$ & $2^{19}$ & 32 MiB \\
Two-candidate table & $2^{29}$ & $2^{11}$ & $2^{20}$ & 64 MiB \\
\bottomrule
\end{tabular}
\caption{Equal-coefficient examples counting 64-byte raw shelter payloads only;
addresses and validity metadata are additional. The second
row assumes $D\simeq2N$ only for sizing; actual hashing parameters must be chosen.}
\end{table}

Both rows have 512 batches per epoch, or approximately 52.43 seconds for one
replica receiving all 10,000 logical requests/s. With $K$ equally loaded
independent replicas, each epoch lasts approximately $52.43K$ seconds; aggregate
amortized rebuild work is unchanged, while state replication increases.
Main-array dummy capacity and payload/share
encoding may increase these raw sizes. Deamortizing rebuilding introduces
additional memory and scheduling costs; the table gives no tail-latency guarantee.

\subsection{Crossover against an efficient direct batch}

For a local approximation $F(D,b)\simeq c_FD$, set
\[
\alpha=c_H/c_F,\qquad \beta=c_R/c_F.
\]
Ignoring additive and recursive differences, the optimally tuned ratio is
\begin{equation}
\frac{C_{\rm sqrt,batch}^*}{F(D,b)}
=2\sqrt{\alpha\beta\,b/D}.
\label{eq:crossover}
\end{equation}
For both tabled layouts $D/b=2^{18}$, hence
\[
\frac{C_{\rm sqrt,batch}^*}{F(D,b)}=\frac{\sqrt{\alpha\beta}}{256}.
\]
When the stationary optimum is feasible, the simplified model favors square-root
ORAM when $\alpha\beta<65{,}536$.
At the untuned choice $S=\sqrt{Db}$ the ratio is instead
$(\alpha+\beta)/512$. With equal constants the modeled ratio is $1/256$;
\emph{this is not an end-to-end speedup prediction}.

The practical question is whether the private shelter join, rebuild, recursion,
and online physical accesses fit within this headroom when compared with a
well-implemented direct gather. A direct scan may have excellent memory
locality and simpler preprocessing. The model establishes a plausible research
direction, not a winner. Benchmarking only a naive $Db$ scan would not answer
the user's question.

\section{Recursive position maps and the large-base proposal}\label{sec:recursion}

Let ORAM level 0 contain $n_0=D$ records. Pack $F_j$ position labels into a record
at the next level, so
\[
n_{j+1}=\left\lceil n_j/F_j\right\rceil.
\]
If each label uses one field element, a level-$(j+1)$ record carries approximately
$F_j$ such elements. Fewer records therefore do not imply proportionally fewer
payload words. A packed block also requires private subblock selection. In the
lookup-only square-root construction, stored parent position labels remain
fixed within the parent epoch and are replaced during bulk parent reconstruction.
The recursive map's own shelter/access metadata still changes. Tree ORAM may
instead update a retrieved label on each access.

If levels $0,\ldots,L$ are ORAMs and the position map of level $L$ is flat, the
flat map has $n_L$ scalar labels. Stopping immediately at level 0 leaves a
$D$-label map. One cannot count a $D/F_0$-entry base without introducing and
paying for the additional packed ORAM level.

A useful batch accounting template is
\begin{align}
C_{\rm recursive}={}&\sum_{j=0}^{L}
\left(J_j(S_j,b)+\frac{b}{S_j}R_j(n_j,w_j)\right)
+F_{\rm base}(n_L,b)\nonumber\\
&+\sum_{j=0}^{L-1}\operatorname{SubblockSelect}_j(b,F_j)
+\operatorname{GroupingAndRouting}.
\label{eq:recursive}
\end{align}
Here using $b$ at every level is a padded schedule; a more refined construction
may exploit private grouping but must not reveal the number of distinct blocks.
Parent rebuilds can force descendant reconstruction. Those costs must be counted
once, at their actual frequency, rather than assigning independent optimal
periods to all levels. In the one-level model, $R_{\rm build}$ may aggregate all
work triggered by that rebuild; in the recursive sum, do not count that same
descendant work again. Likewise, let $J_j$ be level-local and assign cross-level
grouping/routing to one term. Bulk rebuilds are essential; issuing a separate recursive
write per parent entry can destroy the intended savings.

\subsection{Why larger fanout can help, and where it stops helping}

\status{Derived coefficient model}
For $n$ labels packed $F$ per record, suppose shelter work per record is
$\gamma+\mu F$ and rebuild work per record is $\rho+\nu F$.
Optimizing that added level gives
\begin{equation}
2\sqrt{bn\left(\frac{\gamma\rho}{F}
+\gamma\nu+\mu\rho+\mu\nu F\right)}.
\label{eq:fanout}
\end{equation}
If comparison/control dominates shelter work and payload dominates rebuilding,
the cross term can dominate and the cost is approximately insensitive to $F$.
Larger fanout may then reduce recursive depth and base size cheaply. Once both
phases are payload-dominated, the $\mu\nu F$ term penalizes excessive fanout.
Subblock extraction can impose an additional penalty absent from this simplified
formula.

Choose the flat cutoff by comparing its \emph{batched} gather with the complete
cost of one more level plus its smaller base. A large base is plausible because
batched gather is much cheaper than $b$ independent scans. It is not justified
by a universal asymptotic cutoff. Also, a shared epoch size eventually exceeds
the capacity of a small recursive level: stop recursion, explicitly pad that
level, or design a valid nested schedule.

\section{Tree ORAM, eviction, and sparse-key maps}\label{sec:trees}

\subsection{Circuit ORAM and the depth of eviction}

Circuit ORAM is a natural small-circuit reference \cite{circuit2015}. The
concern is that its eviction planning contains dependent scans and metadata
comparisons that are relatively expensive in arithmetic MPC. Nevertheless,
sequential-looking pseudocode does not imply an inherent $\Theta(\log D)$
interactive depth for every eviction implementation.

Chan, Chung, and Shi explicitly parallelize Circuit ORAM's planning and eviction
in the full version of \emph{On the Depth of Oblivious Parallel RAM}, Section 7.1
\cite{ccs2017}. For a path with $L$ levels they obtain $O(\log L)$ word depth and
$O(L\log L)$ work while preserving the sequential eviction outcome. This is the
most directly relevant literature answer to the eviction-depth question.
It does not mean $\lceil\log_2 28\rceil=5$ SPDZ rounds: a word comparison, lookup,
or conditional move must itself be compiled, and paths interact through buckets
and stash state.

Titanium's optimized Circuit ORAM eviction reduces payload conditional
operations using combined pick/drop swaps \cite{titanium2022}. This addresses a
different part of the cost from parallelizing the planning scan. A useful
implementation should measure metadata planning and wide-record movement
separately, then investigate whether their optimizations compose.

\subsection{Arithmetic-friendly reformulations to test}

\status{Implementation ideas, not measured improvements}
For a public eviction path, compute each block's eligibility for all relevant
depths using the bits of its secret leaf label and public path bits. Shared
prefix-match bits can be formed by a parallel prefix product rather than
repeatedly comparing full integers. Keep already-needed metadata in bit form
where this avoids repeated arithmetic-to-bit conversion.

Maxima, first/last eligible positions, and selected-candidate propagation can
often be written as associative scans on tagged summaries. The combination
operator must be associative with explicit tie-breaking and validity rules;
merely calling a loop a prefix scan does not establish this. The published
parallel Circuit ORAM analysis is the correct starting point. A resulting
design still needs a stash-overflow argument and proof that its eviction
schedule preserves the ORAM invariant.

Appendix~\ref{app:eviction} retains more explicit proposed scan formulations
from the discussion. Eviction's bucket-depth/target comparisons involve roughly
five-bit control words at the target scale, rather than 256-bit map keys. Leaf
compatibility uses the longer leaf labels. Storing every leaf bit as a separate
arithmetic share can substantially expand metadata; compact storage and free
bit access cannot both be assumed.

\emph{Efficient Mixed-Mode Oblivious RAMs} also proposes parallel bucket-merge
eviction \cite{mixedmode2025}. It trades dependence depth against larger buckets
and merge circuits. Its specialized three-party design is a candidate to study,
not an established malicious ten-party replacement.

Ring ORAM is worth comparing because it shifts work into bucket reshuffling and
scheduled eviction \cite{ringoram}. However, its client/server bandwidth savings
do not directly imply a smaller malicious MPC circuit. Private slot selection,
validity bookkeeping, reshuffling, and authentication may dominate in this
setting. No verified alternative was found that simultaneously guarantees
Circuit-ORAM-like concrete cost and a superior ten-party arithmetic-MPC
implementation for this workload.

\subsection{What a batch of paths can save}

If $b$ independently uniform leaf paths in a binary tree are fetched together,
the expected number of distinct nodes at depth $d$ is
\begin{equation}
2^d\left(1-(1-2^{-d})^b\right).
\label{eq:path-union}
\end{equation}
Summing over depths gives approximately
$b(\log_2(D/b)+O(1))$ nodes rather than $b\log_2 D$, up to the relation between
leaf count and array capacity. This is an occupancy calculation, not a complete
batched ORAM algorithm. Shared upper nodes offer limited asymptotic savings at
$D=2^{28},b=2^{10}$: roughly the lower 18 levels remain distinct per query.
Repeated logical requests, adaptive position-map lookups, and conflicting
evictions need additional coordination. Public paths must have the distribution
required by the ORAM security proof.
For height 28 and 1,024 independent paths, Equation~\eqref{eq:path-union} gives
19,568.31 expected distinct buckets versus 29,696 separate visits, a 34.10\%
reduction in this count alone. Eviction paths have their own schedule and must
not automatically be assigned the same random-path model.

Jarecki and Wei study three-party Circuit ORAM with fast batch retrieval
\cite{jarecki2018}; FABLE studies efficient batched private lookup in a two-party
setting \cite{fable2025}. Both are relevant design references, but their concrete
protocol costs do not transfer directly to ten malicious parties.
Jarecki--Wei's stated fast-batch regime is $b=O(\log D)$ and tolerates one
semi-honest corruption among three parties; using batch 1,024 requires examining
its schedule and stash constraints. FABLE is semi-honest two-party computation
with a table held by one party, and its server still performs linear table work.
Its low communication should not be interpreted as sublinear total computation.

\subsection{Sparse 256-bit keys: cuckoo table versus an oblivious search tree}

\textbf{Cuckoo representation.} Compute two candidate indices under suitable
hash functions/PRFs inside MPC and read both table locations obliviously. Compare
the stored full keys, then select the value or a missing-key result. Charge two
dense accesses, hash-table slack, PRF evaluation, collision handling, any hashing
stash, and exact key equality. ``Two ORAM accesses'' describes the ordinary
two-candidate lookup only; it omits any overflow structure. Initialization and
failure parameters must be selected for $2^{28}$ entries.
The original two-choice cuckoo construction is a useful reference
\cite{cuckoo2004}; one-slot variants require operating slack around twice the
logical capacity, while bucketed variants change the number of candidate
records compared. Both probes can be evaluated in parallel.

The phrase ``MPC-friendly PRF'' is a design requirement, not a selected primitive.
The relevant metric is arithmetic cost and depth over the actual field, with
appropriate security parameters and domain separation for the two candidates.
Public deterministic hashing may be possible under a different access/leakage
argument; it must not be substituted without specifying that argument. A static
dictionary or perfect-hash construction is also a possible future read-only
alternative, but was not analyzed in this discussion.
MiMC and Poseidon are candidate field-native primitive families
\cite{mimc2016,poseidon2021}, not a selected keyed instantiation. Encoding
256-bit inputs, keyed/domain-separated use, range reduction to array indices,
and multiplication depth all need to be specified and costed.

\textbf{Search tree integrated with location information.} EnigMap is relevant
to overlapping search-tree navigation and ORAM position information
\cite{enigmap}. This can avoid treating every tree edge as an independent access
to a fully recursive generic ORAM. Its secure-enclave/client setting differs
from arithmetic MPC. Moreover, its paper's batch requests are processed
sequentially; a large-batch benchmark does not establish the batched MPC
amortization needed here.
The underlying pointer/position-map overlap originates in oblivious data
structures \cite{ods2014}. A literal EnigMap node contains key-bearing pointers
for itself and two children; with 32-byte keys and a 32-byte value this can
already mean 128 bytes before position labels. This is not a lower bound: dense
node identifiers may permit a different encoding. Do not assume the search-tree
node has the same width as the original 64-byte pair.

For this workload a binary search tree entails many dependent comparisons on
long keys unless the comparison representation is redesigned. The choice is
therefore between a small number of hashed dense lookups plus exact equality
and dependent search/navigation with potentially fewer redundant position-map
operations. Measure both the key-processing cost and the memory-access cost;
do not compare only the number of ORAM calls.

\section{Honest-majority arithmetic and degree-4 sharing}\label{sec:backend}

\subsection{Multiplication is not local share multiplication alone}

Let party $i$ hold $f_x(\alpha_i)$ and $f_y(\alpha_i)$, with degree-4 polynomials
and constants $x,y$. Their local product is a point on $f_xf_y$, which has degree
at most eight. Returning that value as a fresh degree-4 share is incorrect.
Repeated local multiplication rapidly exceeds the reconstruction threshold.

One standard route uses correlated sharings of the same fresh random $r$ at
degrees four and eight. Open a properly randomized sharing of $xy+r$, then
subtract $[r]_4$ to obtain a degree-4 result. Another route uses a checked Beaver
triple $(\sh a,\sh b,\sh{ab})$: open
$\delta=x-a$ and $\epsilon=y-b$, then compute
\begin{equation}
\sh{xy}=\sh{ab}+\delta\sh b+\epsilon\sh a+\delta\epsilon.
\label{eq:beaver}
\end{equation}
All output terms remain degree four. The opening and preprocessing subprotocols
must match the active-security model; naive interpolation is not a way to
correct four malicious errors in a degree-8 sharing among ten parties.
In particular, opening the raw degree-8 product polynomial can reveal information
beyond its constant term. Use the backend multiplication primitive.

\subsection{Practical starting points}

MP-SPDZ supports several distinct backends; its name does not imply that every
backend uses conventional dishonest-majority SPDZ \cite{mpspdzdocs}.

\begin{tabularx}{\textwidth}{@{}lY@{}}
\toprule
Executable & Role in the proposed benchmark \\
\midrule
\path{malicious-shamir-party.x} & Malicious prime-field Shamir baseline. \\
\path{sy-shamir-party.x} & Malicious SPDZ-wise Shamir baseline with authentication. \\
\path{shamir-party.x} & Passive baseline; not a substitute for active security. \\
\path{atlas-party.x} & Documented MP-SPDZ implementation is passive, despite the ATLAS paper also giving malicious protocols. \\
\bottomrule
\end{tabularx}

The Shamir options permit $n=10,t=4$; pin a source revision and explicitly set
party count and threshold when benchmarking \cite{shamiroptions}.
Foundational malicious honest-majority references include Lindell--Nof and
Chida et al. \cite{lindellnof,chida2018}. No benchmark in the discussion establishes
which current malicious backend wins on this database circuit.

\subsection{Batching is not automatically packed Shamir sharing}

With ordinary polynomial packing, degree $d$ supports $k$ secret slots and
privacy threshold $t=d-k+1$. Local products require $2d<n$ for ordinary
degree reduction. Hence
\[
k\leq \left\lfloor\frac{n-1}{2}\right\rfloor-t+1.
\]
At $n=10,t=4$, this gives $k\leq1$. Thus a batch of 1,024 requests cannot simply
be packed into a single ordinary degree-4 Shamir polynomial while retaining
four-corruption privacy. Advanced sharing transformations, different sharing
schemes, or weaker thresholds can change the situation. This elementary bound
does not rule out the advanced packed protocols surveyed below.

\section{Equality on degree-4 Shamir sharings}\label{sec:equality}

The output required by an ORAM is usually a \emph{secret} bit
$\sh e$, where $e=1$ iff $x=y$. Revealing this bit can expose hit/miss or duplicate
information. All secure products below return degree-4 sharings through a valid
backend, and all fresh correlations require the backend's active checks.

\subsection{Exact full-field masked equality}

Let $x,y\in\F_p$ and $\ell=\lceil\log_2p\rceil$. Prepare a fresh uniform
$r\in\F_p$ together with degree-4 shares of its \emph{canonical} bits:
\[
r=\sum_{j=0}^{\ell-1}2^jr_j,\quad r_j\in\{0,1\},\quad0\leq r<p.
\]
Open $c=x-y+r\pmod p$. Write the public canonical bits of $c$ as $c_j$ and form
locally
\begin{equation}
\sh{m_j}=(1-c_j)+(2c_j-1)\sh{r_j}.
\label{eq:match}
\end{equation}
Compute
\begin{equation}
\sh e=\prod_{j=0}^{\ell-1}\sh{m_j}
\label{eq:eqtree}
\end{equation}
using a balanced tree of secure multiplications. Each factor is one precisely
when the corresponding bits match. Therefore $e=1$ iff $c=r$, equivalently
$x=y$. The opened $c$ is uniform because $r$ is fresh and uniform. This masked
equality principle appears in Nishide--Ohta, Section 6.4 \cite{nishide2007}; the
balanced-tree implementation and counts used here are a simple specialization.

\begin{table}[htbp]
\centering
\begin{tabular}{@{}lrrr@{}}
\toprule
Input representation & Factors & Secure products & Product layers \\
\midrule
One element of $\F_{2^{61}-1}$ & 61 & 60 & 6 \\
One element of a 128-bit prime field & 128 & 127 & 7 \\
Five independent 61-bit field limbs & 305 & 304 & 9 \\
\bottomrule
\end{tabular}
\caption{Derived online counts after opening the masked differences; mask
production and backend verification are additional. The last row flattens the
five product trees into one tree.}
\end{table}

For 1,024 equalities in the 61-bit field, all first openings can be batched and
the 61,440 secure products occupy six multiplication layers. A layer is not
necessarily one network round trip; the backend may require additional phases.
Fresh masks are consumed by every equality, even when inputs repeat.

\subsection{Canonical masks and their real preprocessing cost}

Uniform independent bits do not automatically encode a uniform canonical field
element. A conceptually simple method samples $\ell$ shared random bits,
computes a secret acceptance predicate $R<p$, opens that predicate, and repeats
on rejection. Accepted masks are uniform on $\F_p$. This work is offline but not
free. For $p=2^{61}-1$, only the all-ones candidate is invalid, with probability
$2^{-61}$; its bits must not be retained while treating its field encoding as
zero. That mismatch would invalidate the equality test.

Random bits can themselves be generated by standard arithmetic preprocessing.
For example, sample a shared random $u$, securely compute and open $u^2$, reject
zero, choose a public square root $q$, and use $(1+u/q)/2$ as a random bit over
an odd prime field. This is only an algebraic description: malicious generation,
consistency, and opening checks remain necessary.

Reusing an additive mask is unsafe: opening $d_1+r$ and $d_2+r$ reveals
$d_1-d_2$. Reusing a mask across recursive levels or duplicate groups is not
justified by batch processing. Mask-bit storage and production should be measured
alongside ordinary multiplication preprocessing.

\subsection{A two-opening power-tuple variant}\label{sec:power}

\status{Derived optimization; not claimed to be a default library protocol or a new result}
The masked bit matches in~\eqref{eq:match} allow equality to be expressed as a
low-domain polynomial. Define the mismatch count
\[
H=\sum_{j=0}^{\ell-1}(1-m_j)\in\{0,\ldots,\ell\}.
\]
For $p>\ell$, let
\begin{equation}
P(T)=\prod_{j=1}^{\ell}\left(1-\frac{T}{j}\right).
\label{eq:indicator-poly}
\end{equation}
Then $P(0)=1$ and $P(h)=0$ for all integers $1\leq h\leq\ell$, so $P(H)=e$.

In addition to the first mask, preprocess an independent uniform $s\in\F_p$
and \emph{separate, checked degree-4 sharings}
\[
\sh s,\sh{s^2},\ldots,\sh{s^\ell}.
\]
Open $u=H+s$. Since $u$ is now public, calculate the public coefficients
$\beta_j(u)$ of
\[
P(u-X)=\sum_{j=0}^{\ell}\beta_j(u)X^j.
\]
Return
\begin{equation}
\sh e=\sum_{j=0}^{\ell}\beta_j(u)\sh{s^j}
=\sh{P(u-s)}=\sh{P(H)},\qquad \sh{s^0}=1.
\label{eq:power-eq}
\end{equation}
The final operation is linear in already prepared degree-4 sharings. It does not
raise a local Shamir share to the $j$th power. The second opening is uniform,
even conditioned on the first transcript, because $s$ is fresh and independent.

This gives two dependent online opening phases and no online secret products
\emph{inside the equality computation}, after correlations are ready. It shifts
$O(\ell)$ products and correlated storage offline; simple sequential power
generation uses $\ell-1$ products, plus verification. Naive public coefficient
shifting takes $O(\ell^2)$ local operations per equality. The first canonical-bit
mask is still needed, and both types of correlations are single-use.

The tradeoff is lower online depth against preprocessing, memory, and public
arithmetic. A continuously running service must replenish the power tuples, so
the improvement in end-to-end throughput is unresolved. Related constant-round
arithmetic techniques are classical \cite{damgard2006}; the exact formulation
above is recorded as a derivation from the discussion rather than an attributed
measured implementation.

\subsection{Bounded integer equality and field-size slack}

Suppose $x,y\in[0,2^k)$, so $d=x-y$ satisfies $|d|<2^k$. A bounded-integer
masking variant opens
\begin{equation}
c=d+2^k+r_{\rm low}+2^kr_{\rm high},
\label{eq:bounded-eq}
\end{equation}
where $r_{\rm low}$ has $k$ shared random bits and $r_{\rm high}$ supplies
statistical hiding. Choose the field and mask ranges so there is no modular
wraparound. The low $k$ bits of $c$ equal those of $r_{\rm low}$ iff
$d\equiv0\pmod{2^k}$, which under the bound means $d=0$.

This allows a $k$-factor equality tree: $k-1$ products and
$\lceil\log_2k\rceil$ layers after the opening. For a 28-bit address this is
27 products and five layers. However, statistical masking needs approximately
$k+\sigma$ bits plus appropriate headroom in the field. A 61-bit prime cannot
simultaneously provide 28-bit operands and 64 bits of integer-masking slack.
The exact full-field protocol avoids this particular slack requirement by using
all canonical field bits; the active verification soundness is a separate issue.
For a concrete sufficient condition, choose independent uniform $k$-bit
$r_{\rm low}$ and $\sigma$-bit $r_{\rm high}$ and take
$p>2^{k+\sigma}+2^{k+1}$. Relative to a transcript simulated with $d=0$, the
per-test statistical privacy distance is below $2^{-\sigma}$. These errors
compose across tests; correctness under the stated bounds remains exact.

MP-SPDZ's integer equality code contains masking-based paths associated with
Catrina--de Hoogh \cite{catrina2010,mpspdzeq}. Other paths, including a known-prime
path, may use different bit decompositions. Therefore ``28-bit equality costs
27 multiplications'' is a count for the explicitly described construction, not
a claim about every invocation of a high-level equality operator. Malicious
inputs must satisfy the assumed range, either by representation or checking.

\subsection{Full 256-bit keys}

A 256-bit key must be encoded injectively. In $\F_{2^{61}-1}$, four 60-bit limbs
and one 16-bit limb suffice. Independently mask each limb and combine all match
factors. The generic full-field route uses $5\cdot61=305$ factors, 304 products,
and nine product layers. The five initial openings can share one communication
phase. This intentionally simple count does not optimize the shorter last limb.

Alternatively, sum the five mismatch counts and use the power-tuple
construction with a polynomial of degree 305. It still has two opening phases,
but substantially larger offline tuples. If both keys are permanently bit-shared,
each bit match is $1-x_j-y_j+2x_jy_j$; direct equality uses $2k-1$ products over
$k$ bits, with $1+\lceil\log_2k\rceil$ layers. If one operand is public, the
per-bit products disappear. Permanent bit storage may be costly across the
entire database.

Fingerprints can reduce comparison width only by introducing and accounting for
collision probability, or by retaining an exact final check. No fingerprint
parameters or adaptive-query security argument were selected in the discussion.

\subsection{Other zero tests and pitfalls}

Fermat gives the exact formula
\[
e=1-(x-y)^{p-1}.
\]
It requires secure squarings and multiplications with degree reduction, and has
a longer dependency chain than a balanced match tree. It avoids canonical-bit
mask preprocessing, so its full cost can still be competitive on a low-latency
backend. Evaluate a good addition chain rather than assuming a naive binary
exponentiation count.

Multiplying $x-y$ by a fresh nonzero secret random value and opening the product
is an inexpensive \emph{public} zero test. It reveals whether $x=y$ and is
therefore generally unsuitable for internal ORAM membership. Neither this test
nor Fermat exponentiation can be implemented by simply opening or exponentiating
unreduced local Shamir products.

\section{Preprocessing literature: what changes with honest majority}\label{sec:literature}

\subsection{A useful practical reference: ATLAS}

ATLAS \cite{atlas2021} gives honest-majority arithmetic MPC, including malicious
security with abort. Its leading full multiplication communication is about
four field elements per party in the information-theoretic variant and about
two with PRG-based savings; sublinear verification terms remain. Round compression
trades these leading constants for approximately 4.5 and 2.5, respectively.
The paper's synthetic benchmark uses a 61-bit Mersenne field and one AWS
\texttt{c4.large} per party. Table 2 reports 10 million multiplications at depth
20 in 7.892 seconds for nine parties and 8.628 seconds for eleven parties.
These imply roughly 1.27 and 1.16 million multiplications/s per complete party
instance, not per core. The experimental variant is information-theoretic.

The author implementation's timed execution includes offline and online phases
\cite{atlascode}.
\emph{Code provenance:} the inspected author checkout is revision
\path{f7bf4193ca24e780103cd08aa9e8930e272147f8}, whose implementation uses four
software threads. The neighboring party counts are a useful scale reference,
not an exact ten-party result. Nor is the synthetic circuit a database benchmark.
The small benchmark instances cannot hold the target database. Field width,
security checks, circuit depth, memory pressure, and networking prevent a direct
conversion to a modern VM's per-core capacity.

\subsection{Recent honest-majority directions}

The following are relevant after accepting honest majority. None supplies a
verified end-to-end measurement of the proposed database.

\begin{longtable}{@{}>{\raggedright\arraybackslash}p{0.27\textwidth}>{\raggedright\arraybackslash}p{0.65\textwidth}@{}}
\toprule
Work & Relevance and qualification \\
\midrule
\endfirsthead
\toprule Work & Relevance and qualification \\
\midrule
\endhead
Song--Ye, EUROCRYPT 2025 \cite{songye2025} &
Malicious honest-majority MPC using sharing transformation and packed Beaver
correlations. Preparation costs $O(n)$ field elements plus $O(n)$ random OLE
correlations per packed triple. OLE generation is part of the cost, not a free
primitive. Relevant to reducing total communication beyond plain Shamir. \\
\addlinespace
Li--Xing--Yao--Yuan, PKC 2026 \cite{sumcheck2026} &
Sumcheck-based malicious verification. For Shamir, retains $O(M)$ extra
computation and $O(\log M)$ extra communication, replaces certain double
sharings with random sharings, and improves a soundness dependence from
$O(M/|\F|)$ to $O(\log M/|\F|)$. Applicable to the ATLAS line; not a measured
ten-party replacement in this discussion. \\
\addlinespace
Helix, 2025 preprint \cite{helix2025} &
General-$n$ malicious honest-majority inference with compressed checks exploiting
operand reuse and a revised prefix-OR primitive. Reusing a control bit across
record words suggests a possible ORAM connection. Published neural-network
speedups are not generic ORAM or scalar multiplication rates. \\
\addlinespace
Bienstock--Escudero--Polychroniadou, CRYPTO 2026 \cite{crt2026} &
CRT-based sharing gives linear total computation and communication for active
corruptions with $t<(1/2-\varepsilon)n$. Asymptotic constants and the required
gap from one half must be evaluated for ten parties. \\
\addlinespace
Du--Song--Ye, ASIACRYPT 2026 preprint \cite{god2026} &
Honest-majority guaranteed output delivery from pairwise random OLEs, using
detectable packed sharing. Linear online communication in the OLE-hybrid model;
information-theoretic preprocessing, random-oracle, and PCG instantiations have
different total costs. GOD is stronger than the current working abort model. \\
\addlinespace
Song--Ye, 2025 preprint revised 2026 \cite{perfect2025} &
Perfect security with $t<n/3$, using packed bivariate triples. At ten parties this
supports at most three corruptions. Linear online communication is conditioned
on preprocessing; preparing these correlations has an $\Omega(n)$ per-secret
communication lower bound in the stated setting. \\
\bottomrule
\end{longtable}

The field size alone does not specify malicious soundness. For instance, a
check with error proportional to $M/|\F|$ behaves differently from one proportional
to $\log M/|\F|$. At $M\approx10^7$ and a 61-bit field, the former scaling is
roughly $2^{-38}$ before constants or repetitions. This is an illustration of
the need to inspect the actual theorem, not an audited failure probability for
a particular implementation. Extension-field checking or repetition can alter
the cost. Statistical hiding for bounded comparisons is another independent
parameter.

\subsection{How to read multiparty triple-generation claims}

The original question asked specifically about more than two parties. Four
distinctions are indispensable:
\begin{enumerate}
\item Raw Beaver triples do not include the authentication needed by a
conventional full-threshold SPDZ backend.
\item A multiparty raw-correlation construction may have only a two-party
authenticated extension or benchmark.
\item A PCG changes seed communication and local expansion; authenticated setup,
checks, memory, and expansion CPU can remain substantial.
\item A protocol that avoids ordinary authenticated triples by changing online
verification is not a drop-in faster triple generator for an unchanged backend.
\end{enumerate}

Relevant dishonest-majority results are retained in Appendix~\ref{app:dm} for
continuity. Their relevance is secondary once a suitable honest-majority backend
is allowed. Improvements for matrix triples or 1024-bit scalar fields should not
be applied to 61-/128-bit scalar ORAM gates without evidence.

\section{AWS cost model for 10,000 reads per second}\label{sec:cost}

\subsection{What is known and what is missing}

\textbf{Known requirement:} $R=10^4$ logical reads/s, ten parties, service running
continuously, with multiple workers allowed. \textbf{Missing measurement:} the
complete work and traffic per lookup for a chosen construction. Therefore the
honest-majority discussion supports a parameterized estimate, not a final
annual bill.

The earlier dishonest-majority estimate of roughly \$0.3--1 billion/year across
all parties assumed 10,000--30,000 authenticated scalar multiplications per
read and extrapolated conventional LowGear-style preprocessing. Both the
per-read gate budget and its ten-party throughput extrapolation were provisional.
\textbf{Do not reuse that estimate as the honest-majority budget.}
The workload has not been shown to require $G=10{,}000$--$30{,}000$; those numbers
are sensitivity scenarios only.

\subsection{Throughput and leading multiplication traffic}

Let $c$ be average transmitted field words per party per full multiplication,
$w$ the encoded field-word size in bytes, and $G$ the complete multiplication
count per logical lookup. For an ATLAS-inspired sensitivity model choose $c=4$:
\begin{align}
M_{\rm sec}&=RG,\\
B_{\rm party/sec}&\simeq RGcw,\label{eq:bandwidth}\\
B_{\rm all/year}&\simeq nRGcw\cdot31{,}536{,}000.
\end{align}
The coefficient is a leading communication proxy, not a statement that every
Shamir implementation has that cost. Rounding effects, verification, openings,
bit masks, permutation protocols, PRFs, framing, and other work may add traffic.
If all those primitives are charged through an equivalent $G$, avoid counting
their products again separately.

\begin{table}[htbp]
\centering
\begin{tabular}{@{}rrrrr@{}}
\toprule
$G$/read & Multiplications/s & Word bytes & GB/s/party & Gbit/s/party \\
\midrule
10,000 & $10^8$ & 8 & 3.2 & 25.6 \\
10,000 & $10^8$ & 16 & 6.4 & 51.2 \\
30,000 & $3\cdot10^8$ & 8 & 9.6 & 76.8 \\
30,000 & $3\cdot10^8$ & 16 & 19.2 & 153.6 \\
\bottomrule
\end{tabular}
\caption{Derived leading traffic at $c=4$, decimal GB and Gbit. These are
conditional scenarios, not measured lookup costs.}
\end{table}

Using $c=2$ halves this leading term but does not establish a matching compute
speedup or a ready-made implementation. Similarly, arithmetic over 16-byte words
cannot be assigned the measured 61-bit multiplication rate merely by doubling
network bytes.

For perspective only, dividing 100--300 million multiplications/s by a nearby
published rate of 1.2 million/s yields about 84--250 parallel MPC worker groups.
Each group has a participant at every party. These are neither 84--250 CPU cores
nor a validated deployment plan. If every group duplicates full ORAM state,
per-party memory can grow into terabytes. Sharing state between workers requires
a correct concurrent ORAM schedule; simply starting more processes is not free
parallelism.

\subsection{Compute-price snapshot}

The following Linux On-Demand, shared-tenancy prices for US East (N. Virginia)
were retrieved in the discussion from AWS's official offer feed
\cite{awsprices}. The stored feed metadata gives publication
\texttt{2026-10-09T20:31:34Z}, version \texttt{20261009203134}, with the selected
rates effective 1 October 2026. They are a dated reproducibility snapshot.

\begin{table}[htbp]
\centering
\small
\begin{tabular}{@{}lrrrrr@{}}
\toprule
Instance & vCPU & RAM & Network & USD/hour & Ten nodes/year \\
\midrule
\texttt{c7i.8xlarge} & 32 & 64 GiB & 12.5 Gbit/s & 1.42800 & 125,093 \\
\texttt{r7i.16xlarge} & 64 & 512 GiB & 25 Gbit/s & 4.23360 & 370,863 \\
\texttt{r6in.16xlarge} & 64 & 512 GiB & 100 Gbit/s & 5.57856 & 488,682 \\
\bottomrule
\end{tabular}
\caption{Ten nodes means one node per party, 8,760 hours/year. This is an
instance-price calculation, not a claim that ten nodes meet the throughput target.}
\end{table}

For $K$ homogeneous nodes per party at price $p_{\rm hr}$,
\begin{equation}
C_{\rm compute/year}=10Kp_{\rm hr}\cdot8760.
\label{eq:compute-cost}
\end{equation}
Choose $K$ to satisfy measured computation, NIC, memory capacity, and memory
bandwidth requirements simultaneously. Peak advertised network performance is
not a guaranteed sustained MPC rate, and vCPU count is not physical-core count.
Storage, initialization, backups, engineering, spare capacity, and operational
overhead are not priced in this table.

\subsection{Network placement can dominate the annual bill}

AWS states that VPC peering traffic within the same physical availability zone
can be free of peering data-transfer charges, including across accounts
\cite{awspeering}. Availability-zone names can differ between accounts; placement
must refer to the same physical zone. Cross-AZ same-region transfer is charged
at \$0.01/GB for each endpoint under the cited VPC pricing, giving \$0.02 per GB
of wire data when counting both billed endpoints \cite{awsvpc}.

Let $f$ be the fraction of the aggregate transmitted bytes that cross zones.
With decimal GB, the corresponding sensitivity model is
\begin{equation}
C_{\rm transfer/year}
=\frac{nRGcw\cdot31{,}536{,}000}{10^9}
\cdot f\cdot0.02.
\label{eq:transfer-cost}
\end{equation}
Here aggregate transmitted bytes count each network send once; the price factor
accounts for billing at both ends. Do not double-count receivers again.
The displayed dollar amounts assume $10^9$ bytes per pricing unit; confirm that
conversion against the applicable AWS usage record before producing a budget.
If the billed unit is $2^{30}$ bytes, the corresponding estimates are 6.87\%
lower. This unit assumption is separate from the much larger uncertainty in $G$.

\begin{table}[htbp]
\centering
\begin{tabular}{@{}rrr@{}}
\toprule
$G$/read & 8-byte words, all cross-AZ & 16-byte words, all cross-AZ \\
\midrule
10,000 & \$20,183,040/year & \$40,366,080/year \\
30,000 & \$60,549,120/year & \$121,098,240/year \\
\bottomrule
\end{tabular}
\caption{All ten parties, $c=4$, $R=10^4$/s, $f=1$, transfer only. Multiply by
$f$ for the modeled crossing fraction.}
\end{table}

Same-AZ placement removes this particular fee under the cited connectivity
conditions; it does not remove NIC limits, compute costs, or shared failure
domains. Different-region pricing and latency require a separate calculation.
Network topology is therefore a material architectural input, not a small
adjustment to a CPU-only estimate.

\section{Recommended experiments and unresolved decisions}\label{sec:next}

\subsection{Experiments in dependency order}

\begin{enumerate}
\item \textbf{Pin the arithmetic environment.} Select malicious-with-abort or
passive security explicitly, $n=10,t=4$ if retaining the working threshold,
field representation, comparison bounds, statistical security, and source
revision. Run preprocessing and online execution together for steady-state
measurements.
\item \textbf{Measure primitive costs.} Benchmark full-field equality, bounded
address equality, ordering, prefix operations, public/secret selection, private
permutation, and stable merging at actual metadata and payload widths. Include
random-bit masks and, if used, power-tuple generation and coefficient evaluation.
\item \textbf{Implement the fair direct baseline.} Use the same dense-array
layout and batch size as the ORAM candidates. Record permutation preparation,
payload application, local RAM traffic, duplicate behavior, and output routing.
\item \textbf{Prototype one square-root level.} Sweep shelter size around the
coefficient-aware optimum. Measure shelter lookup/maintenance and a complete
rebuild separately, then verify their amortized sum over whole epochs.
\item \textbf{Add recursion selectively.} Sweep fanouts and base cutoffs.
Include packed subblock extraction, duplicate position-map blocks, and forced
descendant rebuilds. Compare each added level with its saved flat-base cost.
\item \textbf{Compare a tree candidate.} Start from Circuit ORAM with the
published parallel eviction formulation. Separate narrow control work from
wide payload movement and test safe batch scheduling. Include stash and bucket
parameters in both the security and memory accounts.
\item \textbf{Measure the complete map lookup.} Add key-to-index processing,
all candidate locations, full-key equality, missing-key behavior, and overflow
handling. Only now determine $G$, bytes/read, and a realistic throughput range.
\item \textbf{Scale workers and price the service.} Determine state per worker,
whether replicas or coordinated shared state are used, batch formation delay,
throughput saturation, rebuild tail latency, and network placement. Replace the
illustrative values in Section~\ref{sec:cost} with measured values.
\end{enumerate}

\subsection{Security and construction questions still open}

The square-root design needs a complete argument for duplicate requests,
cover probes, repeated accesses across batches, dummy exhaustion, and recursive
rebuild coupling. The sorted-shelter optimization also needs a concrete secure
merge with favorable costs at a highly asymmetric list ratio. A read-only
application simplifies record consistency but does not eliminate these issues.

For tree ORAM, establish whether the desired batch schedule preserves the
published stash bound and eviction behavior. Independent path processing can
conflict in shared buckets; delayed eviction can change overflow probabilities.
Word-depth improvements must be translated into the selected arithmetic circuit.

For preprocessing, choose a malicious backend rather than combining an online
rate from one protocol, a triple rate from another, and a verification argument
from a third. If adopting newer packed, PCG, or sumcheck machinery, check its
field, corruption threshold, setup assumptions, composability, and output-release
semantics against the database's public memory-access transcript.

\subsection{Decision log for the next session}

\begin{tabularx}{\textwidth}{@{}p{0.29\textwidth}Y@{}}
\toprule
Decision or question & Current state \\
\midrule
Database scale and payload & Fixed user requirements: $2^{28}$ pairs, 32-byte keys and values. \\
Read throughput & Target 10,000/s across ten parties; parallel instances allowed. \\
Application writes & Excluded from the present performance target. \\
Honest majority & Accepted by user. \\
Active versus passive & Active with abort used as the working interpretation; not finalized. \\
Threshold & Degree-4 Shamir suggests $t=4$ among ten; retain unless changed explicitly. \\
Batch size & 1,024 logical requests as the main example; tune against latency and load. \\
Key map & Cuckoo-style dense conversion versus integrated search tree remains open. \\
ORAM family & Enlarged square-root and parallelized tree candidates remain open. \\
Base cutoff and fanout & Tune with complete batch costs, not asymptotics alone. \\
Equality & Explicit exact protocols available; fastest full-cost choice unmeasured. \\
Preprocessing backend & Honest-majority Shamir baselines first; newer methods worth evaluating. \\
Annual total cost & Unresolved until complete primitive and lookup measurements exist. \\
\bottomrule
\end{tabularx}

\section{Notes for future readers and reproducibility}\label{sec:provenance}

This is a consolidation of a research conversation, with primary references
checked where needed. It preserves proposed constructions as proposals and
conditional arithmetic as conditional arithmetic. It does not claim an original
publication, a finished security proof, a deployed system, or an AWS quote.
Prices, software behavior, and preprint versions may change.

The source is a single self-contained LaTeX file with an embedded bibliography.
Compile with \texttt{latexmk -pdf batched\_mpc\_oram\_report.tex}, or run
\texttt{pdflatex} twice. No external figures, bibliography database, or local
research checkout is required. Primary URLs in the bibliography make later
verification possible without the chat transcript.

To continue efficiently, read Sections~\ref{sec:restart}, \ref{sec:next}, and
\ref{sec:cost} first. The highest-value next artifact is a reproducible benchmark
of batched gather, shelter merge, rebuild, and equality on the chosen ten-party
backend. That benchmark should settle the coefficient ratios in
Equation~\eqref{eq:crossover} before committing to a complete ORAM construction.

\appendix
\section{Dishonest-majority multiparty preprocessing references}\label{app:dm}

This appendix retains the earlier literature question without treating it as
the main deployment assumption.

\begin{description}
\item[Abram--Scholl, PKC 2022 \cite{abramscholl}.]
Ring-LPN-based multiparty authenticated triples for SPDZ over large prime fields,
including a malicious distributed setup. It is a genuine multiparty reference;
party-count-dependent setup and seed costs require attention. No implementation
from this work was used in the discussion.
\item[Multipars, PoPETs 2024 \cite{multipars}.]
Authenticated multiparty preprocessing over $\mathbb Z_{2^k}$ using linearly
homomorphic encryption, with public code. The concrete evaluation discussed was
two-party; the protocol's multiparty applicability does not supply a ten-party
throughput measurement. Rings also change the equality/encoding choices.
\item[MatriGear, IEEE S\&P 2025 \cite{matrigear}.]
Improves authenticated matrix and large-prime-field triple generation through
HE packing. Its reported 3.6-fold scalar improvement concerns a 1024-bit prime;
the 34-fold result concerns $128\times128$ matrix triples. Neither is a speedup
factor for 128-bit scalar ORAM multiplications.
\item[Boyle et al., CRYPTO 2026 \cite{pirauth2026}.]
\emph{Dishonest-Majority Secure Computation via PIR-Authenticated Multiplication
Triples} bypasses ordinary AMT generation through a different authentication
strategy built around raw PCG correlations and authenticated PIR. Multiparty
extensions are relevant, but this is not an unchanged SPDZ backend with a faster
drop-in AMT factory. No ten-party database benchmark was established.
\item[Li--Song, ASIACRYPT 2026 preprint \cite{free2026}.]
\emph{Malicious Security Comes Free in SPDZ} achieves a leading $4n$ words per
multiplication in the maximum-corruption setting using tensor-product PCG
correlations and recursive checking. Additional local computation remains.
Its verification architecture differs from authenticating every wire in
conventional SPDZ.
\end{description}

Recent PCG papers on arbitrary finite fields, $\mathbb Z/p^k\mathbb Z$, sparse
LPN, and Boolean triples were also examined \cite{pcgfield2025,pcgring2025,
pcfsparse2025,pcgf22025}. Their multiparty raw-correlation results must be
distinguished from authenticated variants that may be only two-party. Similarly,
two-party batch-authenticated PCG speedups \cite{batchauth2026} are not evidence
of multiparty authenticated throughput. Parameter revisions following attacks
on assumptions must be respected; old slide throughput should not override a
revised paper's parameters.

\section{Small reproducible numerical calculations}\label{app:calculator}

The following Python reproduces the report's conditional arithmetic; it is not
an MPC benchmark. Network units are decimal, while storage GiB/MiB is binary.

\begin{lstlisting}[language=Python]
from math import sqrt, log1p, expm1
N, Q, rate, parties = 2**28, 1024, 10_000, 10
seconds = 365 * 24 * 3600
for D, b in [(N, Q), (2*N, 2*Q)]:
    shelter = sqrt(D*b)
    print("S, batches/epoch, seconds/epoch:",
          shelter, shelter/b, shelter/b*Q/rate)
    print("ideal optimized/direct ratio:", 2*sqrt(b/D))
for gates in [10_000, 30_000]:
    for word_bytes in [8, 16]:
        sent = rate * gates * 4 * word_bytes  # per party/sec
        cross_az_cost = sent * parties * seconds / 1e9 * 0.02
        print(gates, word_bytes, sent/1e9, sent*8/1e9,
              cross_az_cost)
for name, hourly in [("c7i.8xlarge", 1.428),
                     ("r7i.16xlarge", 4.2336),
                     ("r6in.16xlarge", 5.57856)]:
    print(name, hourly * 8760 * parties)
# Expected union of 1024 paths, depths 0 through 28 inclusive:
nodes = 1 + sum(2**d * -expm1(Q*log1p(-2**(-d)))
                for d in range(1, 29))
print("expected distinct buckets:", nodes)
\end{lstlisting}

\section{Proposed scan circuits for Circuit ORAM eviction}\label{app:eviction}

\status{Derived circuit reformulations; not a complete batched ORAM proof}
This appendix retains concrete ideas from the discussion. The published
Chan--Chung--Shi construction remains the established reference. The scans below
require implementation-level checking against the chosen Circuit ORAM semantics,
including stash handling, tie-breaking, and bucket-local selection. They do not
establish a stash bound for delayed batch eviction.

\subsection{Eligibility from a public path}

For public eviction-path bits $p_1,\ldots,p_h$ and shared leaf bits
$x_1,\ldots,x_h$, set $m_j=x_j$ if $p_j=1$ and $m_j=1-x_j$ otherwise. Then
\[
E_d=\prod_{j=1}^d m_j,\qquad g=\sum_{d=1}^h E_d
\]
give the allowed-depth indicators and common-prefix length. A work-efficient
prefix-product scan has $O(h)$ multiplication work and $O(\log h)$ depth.
Authenticated bit availability is an assumption; conversion from compact field
labels must be charged if needed.

\subsection{A proposed target planner}

Number path locations $0,\ldots,L-1$, including a stash location only with a
consistent renumbering of all permissible destinations. Let $g_i$ be the maximum
permissible destination among location $i$'s real blocks, or $-1$ if none, and
let $e_i$ indicate an originally empty slot. Ties favor the shallower source.
Define
\begin{align*}
G_i&=\max_{0\leq j\leq i}g_j,\qquad G_{-1}=-1,\\
r_i&=[g_i>G_{i-1}]\land[g_i\geq0],\\
c_d&=[G_{d-1}\geq d]\quad(d\geq1).
\end{align*}
Here brackets denote secret predicates, not public revelations. A source can
win the greedy best-ancestor choice only if it sets a new prefix record.

Compute backward from $u_{L-1}=0$:
\[
u_i=c_{i+1}\land(e_{i+1}\lor u_{i+1}),\qquad
a_i=r_i\land u_i,\qquad v_i=e_i\lor a_i.
\]
The intended meaning is that $a_i$ selects exporting sources and $v_i$ indicates
a vacancy after their export. For the target scan use a distinguished absent
value $\bot$, represented by a separate presence bit. Put $t_{L-1}=\bot$ and,
with $d=i+1$, compute
\[
t_i=\begin{cases}
\bot,&\neg c_d,\\
d,&c_d\land r_d\land v_d,\\
\bot,&c_d\land r_d\land\neg v_d,\\
t_{i+1},&c_d\land\neg r_d\land(t_{i+1}\ne\bot),\\
d,&c_d\land\neg r_d\land(t_{i+1}=\bot)\land v_d,\\
\bot,&\text{otherwise}.
\end{cases}
\]
Finally $\operatorname{target}_i=t_i$ for exporters and $\bot$ otherwise.

\textbf{Reasoning to retain.} Let the record sources be
$q_0<q_1<\cdots<q_m$, with $q_{m+1}=L$. Source $q_j$ is the best eligible
ancestor on the interval
\[
I_j=(q_j,\min\{g_{q_j},q_{j+1},L-1\}].
\]
A vacancy comes from an original empty slot in this interval, or from the next
record source at its endpoint if that source itself exports. Coverage gaps
prevent propagation. The $u$ recurrence captures availability, and the $t$
recurrence chooses the deepest eligible vacancy, resetting at a gap or the next
record source. This is the intended correspondence to greedy sequential eviction;
it should be formalized against the exact bucket/stash invariants before use.

\subsection{Associative summaries for the backward passes}

The $u$ step is an affine Boolean map over the arithmetic field:
\[
f_d(x)=A_d+B_dx,\qquad A_d=c_de_d,\quad B_d=c_d-A_d.
\]
Composition is associative:
\[
(A,B)\circ(A',B')=(A+BA',BB').
\]
For the target step, use a reset bit $R$ and optional candidate $C$:
\[
F_{R,C}(x)=\begin{cases}
C,&R=1,\\
x,&R=0\text{ and }x\ne\bot,\\
C,&R=0\text{ and }x=\bot.
\end{cases}
\]
At destination $d$, take $R_d=\neg c_d\lor r_d$ and let $C_d=d$ if
$c_d\land v_d$, otherwise $\bot$. Composition yields
\[
F_{R,C}\circ F_{S,D}=F_{R\lor S,\ C'},\qquad
C'=\begin{cases}
C,&R=1,\\
D,&R=0\text{ and }D\ne\bot,\\
C,&\text{otherwise}.
\end{cases}
\]
These summaries support $O(L)$ scan-combination work and $O(\log L)$ depth.
Prefix maxima, small comparisons, and bucket candidate selection remain
additional work. The scan's constant depth is not one multiplication layer per
doubling step for the entire eviction.

\subsection{Parallel payload propagation}

Once targets are known, preselect each bucket's original export block $D_i$ and
form optional records $(a_i,D_i,\operatorname{target}_i)$. An exclusive prefix
scan selecting the rightmost valid record gives the most recent exporting source
strictly above each bucket. Accept its payload at bucket $i$ precisely when its
target is $i$.

The intended invariant is that the greedy single-carried-block eviction has
nonoverlapping source-to-destination intervals. An earlier block already deposited
has a smaller target and is rejected by the target equality. This avoids an
$L$-by-$L$ payload routing matrix. With bucket capacity $Z$ and payload width
$w$ field words, the proposed work is $O(LZw+Lw)$ and payload scan depth
$O(\log L)$, plus control and local selection. This is a circuit-level proposal,
not a claim of measured improvement. Compare its constants with Titanium-style
sequential swaps rather than multiplying two independent advertised speedups.

\clearpage
\phantomsection
\addcontentsline{toc}{section}{References}
\begin{thebibliography}{99}
\raggedright

\bibitem{laud2015}
Peeter Laud.
\emph{Parallel Oblivious Array Access for Secure Multiparty Computation and Privacy-Preserving Minimum Spanning Trees}.
Proceedings on Privacy Enhancing Technologies, 2015(2):188--205, 2015.
\url{https://petsymposium.org/popets/2015/popets-2015-0011.pdf}.

\bibitem{falk2021}
Brett Hemenway Falk and Rafail Ostrovsky.
\emph{Secure Merge with $O(n\log\log n)$ Secure Operations}.
ITC 2021, LIPIcs 199, article 7.
\url{https://drops.dagstuhl.de/entities/document/10.4230/LIPIcs.ITC.2021.7}.

\bibitem{blunk2025}
Mark Blunk, Paul Bunn, Samuel Dittmer, Steve Lu, and Rafail Ostrovsky.
\emph{Linear-Time Secure Merge in $O(\log\log n)$ Rounds}.
ITC 2025, article 7.
\url{https://drops.dagstuhl.de/entities/document/10.4230/LIPIcs.ITC.2025.7}.

\bibitem{logstar}
Suvradip Chakraborty, Stanislav Peceny, Srinivasan Raghuraman, and Peter Rindal.
\emph{Logstar: Efficient Linear* Time Secure Merge}.
Cryptology ePrint Archive, Paper 2024/159, revised September 2026.
\url{https://eprint.iacr.org/2024/159}.

\bibitem{shuffle2024}
Jiacheng Gao, Yuan Zhang, Sheng Zhong, and Changyu Dong.
\emph{Free Linear Online Phase for Secure Multiparty Shuffle}.
Cryptology ePrint Archive, Paper 2024/1936, revised September 2026;
ASIACRYPT 2026 version.
\url{https://eprint.iacr.org/2024/1936}.

\bibitem{zahur2016}
Samee Zahur, Xiao Wang, Mariana Raykova, Adri\`a Gasc\'on, Jack Doerner, David Evans, and Jonathan Katz.
\emph{Revisiting Square-Root ORAM: Efficient Random Access in Multi-Party Computation}.
IEEE Symposium on Security and Privacy, 2016.
\url{https://www.cs.umd.edu/~jkatz/papers/sqoram.pdf}.

\bibitem{circuit2015}
Xiao Wang, T-H.~Hubert Chan, and Elaine Shi.
\emph{Circuit ORAM: On Tightness of the Goldreich-Ostrovsky Lower Bound}.
ACM CCS 2015, pp.~850--861; full version, Cryptology ePrint Archive, Paper 2014/672.
\url{https://eprint.iacr.org/2014/672}.

\bibitem{ccs2017}
T-H.~Hubert Chan, Kai-Min Chung, and Elaine Shi.
\emph{On the Depth of Oblivious Parallel RAM}.
ASIACRYPT 2017; full version, Cryptology ePrint Archive, Paper 2017/861.
\url{https://eprint.iacr.org/2017/861}.

\bibitem{titanium2022}
Weikeng Chen, Thang Hoang, Jorge Guajardo, and Attila A.~Yavuz.
\emph{Titanium: A Metadata-Hiding File-Sharing System with Malicious Security}.
NDSS 2022; full version, Cryptology ePrint Archive, Paper 2022/051.
\url{https://eprint.iacr.org/2022/051}.

\bibitem{ringoram}
Ling Ren, Christopher Fletcher, Albert Kwon, Emil Stefanov, Elaine Shi, Marten van Dijk, and Srinivas Devadas.
\emph{Constants Count: Practical Improvements to Oblivious RAM}.
USENIX Security, 2015.
\url{https://www.usenix.org/conference/usenixsecurity15/technical-sessions/presentation/ren-ling}.

\bibitem{jarecki2018}
Stanislaw Jarecki and Boyang Wei.
\emph{3PC ORAM with Low Latency, Low Bandwidth, and Fast Batch Retrieval}.
ACNS 2018, pp.~360--378.
\url{https://eprint.iacr.org/2018/347}.

\bibitem{fable2025}
Zhengyuan Su, Qi Pang, Simon Beyzerov, and Wenting Zheng.
\emph{FABLE: Batched Evaluation on Confidential Lookup Tables in 2PC}.
USENIX Security, 2025.
\url{https://www.usenix.org/conference/usenixsecurity25/presentation/su-zhengyuan}.

\bibitem{enigmap}
Afonso Tinoco, Sixiang Gao, and Elaine Shi.
\emph{EnigMap: External-Memory Oblivious Map for Secure Enclaves}.
USENIX Security, 2023, pp.~4033--4050.
\url{https://www.usenix.org/conference/usenixsecurity23/presentation/tinoco}.

\bibitem{mpspdzdocs}
MP-SPDZ contributors.
\emph{MP-SPDZ Documentation: Getting Started and Protocol Overview}.
Maintained software documentation; accessed 13 October 2026.
\url{https://mp-spdz.readthedocs.io/en/latest/readme.html}.

\bibitem{shamiroptions}
MP-SPDZ contributors.
\emph{MP-SPDZ: Shamir Protocol Options}.
Maintained source code; accessed 13 October 2026.
\url{https://github.com/data61/MP-SPDZ/blob/master/Protocols/ShamirOptions.cpp}.

\bibitem{lindellnof}
Yehuda Lindell and Ariel Nof.
\emph{A Framework for Constructing Fast MPC over Arithmetic Circuits with Malicious Adversaries and an Honest-Majority}.
ACM CCS, 2017.
\url{https://doi.org/10.1145/3133956.3133999}.

\bibitem{chida2018}
Koji Chida, Daniel Genkin, Koki Hamada, Dai Ikarashi, Ryo Kikuchi, Yehuda Lindell, and Ariel Nof.
\emph{Fast Large-Scale Honest-Majority MPC for Malicious Adversaries}.
CRYPTO 2018, pp.~34--64.
\url{https://eprint.iacr.org/2018/570}.

\bibitem{nishide2007}
Takashi Nishide and Kazuo Ohta.
\emph{Multiparty Computation for Interval, Equality, and Comparison Without Bit-Decomposition Protocol}.
PKC 2007, pp.~343--360.
\url{https://doi.org/10.1007/978-3-540-71677-8_23}.

\bibitem{damgard2006}
Ivan Damg\aa rd et al.
\emph{Unconditionally Secure Constant-Rounds Multi-party Computation for Equality, Comparison, Bits and Exponentiation}.
TCC 2006, pp.~285--304.
\url{https://doi.org/10.1007/11681878_15}.

\bibitem{catrina2010}
Octavian Catrina and Sebastiaan de Hoogh.
\emph{Improved Primitives for Secure Multiparty Integer Computation}.
SCN 2010, pp.~182--199.
\url{https://doi.org/10.1007/978-3-642-15317-4_13}.

\bibitem{mpspdzeq}
MP-SPDZ contributors.
\emph{MP-SPDZ Compiler: Equality-to-Zero and Nonlinear Arithmetic Implementations}.
Maintained source code, \texttt{Compiler/floatingpoint.py} and
\texttt{Compiler/non\_linear.py}; accessed 13 October 2026.
\url{https://github.com/data61/MP-SPDZ/blob/master/Compiler/floatingpoint.py};
\url{https://github.com/data61/MP-SPDZ/blob/master/Compiler/non_linear.py}.

\bibitem{atlas2021}
Vipul Goyal, Hanjun Li, Rafail Ostrovsky, Antigoni Polychroniadou, and Yifan Song.
\emph{ATLAS: Efficient and Scalable MPC in the Honest Majority Setting}.
CRYPTO 2021; full version, Cryptology ePrint Archive, Paper 2021/833.
\url{https://eprint.iacr.org/2021/833}.

\bibitem{atlascode}
ATLAS authors.
\emph{ATLAS: Reference Implementation and README}.
Maintained software repository; accessed 13 October 2026.
\url{https://github.com/hanjunli/ATLAS}.

\bibitem{songye2025}
Yifan Song and Xiaxi Ye.
\emph{Honest Majority MPC with $\tilde{O}(|C|)$ Communication in Minicrypt}.
EUROCRYPT 2025; Cryptology ePrint Archive, Paper 2025/283.
\url{https://eprint.iacr.org/2025/283}.

\bibitem{sumcheck2026}
Zhe Li, Chaoping Xing, Yizhou Yao, and Chen Yuan.
\emph{On the Power of Sumcheck in Secure Multiparty Computation}.
PKC 2026; full version, Cryptology ePrint Archive, Paper 2025/177.
\url{https://eprint.iacr.org/2025/177}.

\bibitem{helix2025}
Yansong Zhang, Xiaojun Chen, Qinghui Zhang, Xudong Chen, and Ye Dong.
\emph{Helix: Scalable Multi-Party Machine Learning Inference against Malicious Adversaries}.
Preprint, Cryptology ePrint Archive, Paper 2025/347, 2025.
\url{https://eprint.iacr.org/2025/347}.

\bibitem{crt2026}
Alexander Bienstock, Daniel Escudero, and Antigoni Polychroniadou.
\emph{Actively Secure MPC with $O(|C|)$ Computation and Communication via CRT}.
CRYPTO 2026; Cryptology ePrint Archive, Paper 2026/1270.
\url{https://eprint.iacr.org/2026/1270}.

\bibitem{god2026}
Zonghang Du, Yifan Song, and Xiaxi Ye.
\emph{Honest Majority MPC Achieving GOD with $O(|C|)$ Communication From Pairwise Random OLE Correlations}.
ASIACRYPT 2026 preprint, Cryptology ePrint Archive, Paper 2026/1948.
\url{https://eprint.iacr.org/2026/1948}.

\bibitem{perfect2025}
Yifan Song and Xiaxi Ye.
\emph{Scalable Secure Multiparty Computation with Perfect Security from Preprocessing}.
Preprint, Cryptology ePrint Archive, Paper 2025/1406, revised 2026.
\url{https://eprint.iacr.org/2025/1406}.

\bibitem{awsprices}
Amazon Web Services.
\emph{Amazon EC2 On-Demand Price List: US East (N.~Virginia)}.
Maintained price feed; publication snapshot 9 October 2026,
20:31:34 UTC, version \texttt{20261009203134}; accessed 13 October 2026.
\url{https://pricing.us-east-1.amazonaws.com/offers/v1.0/aws/AmazonEC2/current/us-east-1/index.csv}.

\bibitem{awspeering}
Amazon Web Services.
\emph{What Is VPC Peering?}
Amazon Virtual Private Cloud User Guide, including pricing for a VPC peering connection.
Maintained documentation; accessed 13 October 2026.
\url{https://docs.aws.amazon.com/vpc/latest/peering/what-is-vpc-peering.html}.

\bibitem{awsvpc}
Amazon Web Services.
\emph{Amazon VPC Pricing}.
Maintained pricing documentation; accessed 13 October 2026.
\url{https://aws.amazon.com/vpc/pricing/}.

\bibitem{abramscholl}
Damiano Abram and Peter Scholl.
\emph{Low-Communication Multiparty Triple Generation for SPDZ from Ring-LPN}.
PKC 2022, pp.~221--251.
\url{https://iacr.org/archive/pkc2022/131770119/131770119.pdf}.

\bibitem{multipars}
Sebastian Hasler, Pascal Reisert, Marc Rivinius, and Ralf K\"usters.
\emph{Multipars: Reduced-Communication MPC over $\mathbb{Z}_{2^k}$}.
Proceedings on Privacy Enhancing Technologies, 2024(2):5--28, 2024.
\url{https://petsymposium.org/popets/2024/popets-2024-0038.php}.

\bibitem{matrigear}
HyunHo Cha, Intak Hwang, Seonhong Min, Jinyeong Seo, and Yongsoo Song.
\emph{MatriGear: Accelerating Authenticated Matrix Triple Generation with Scalable Prime Fields via Optimized HE Packing}.
IEEE Symposium on Security and Privacy, 2025, pp.~2453--2471.
\url{https://doi.org/10.1109/SP61157.2025.00174}.

\bibitem{pirauth2026}
Elette Boyle, Niv Gilboa, Matan Hamilis, Yuval Ishai, and Ariel Nof.
\emph{Dishonest-Majority Secure Computation via PIR-Authenticated Multiplication Triples}.
CRYPTO 2026; Cryptology ePrint Archive, Paper 2026/257.
\url{https://eprint.iacr.org/2026/257}.

\bibitem{free2026}
Junru Li and Yifan Song.
\emph{Malicious Security Comes Free in SPDZ}.
ASIACRYPT 2026 preprint, Cryptology ePrint Archive, Paper 2026/283.
\url{https://eprint.iacr.org/2026/283}.

\bibitem{pcgfield2025}
Zhe Li, Chaoping Xing, Yizhou Yao, and Chen Yuan.
\emph{Efficient Pseudorandom Correlation Generators for Any Finite Field}.
EUROCRYPT 2025; Cryptology ePrint Archive, Paper 2025/169.
\url{https://eprint.iacr.org/2025/169}.

\bibitem{pcgring2025}
Zhe Li, Chaoping Xing, Yizhou Yao, and Chen Yuan.
\emph{Efficient Pseudorandom Correlation Generators over $\mathbb{Z}/p^k\mathbb{Z}$}.
CRYPTO 2025; Cryptology ePrint Archive, Paper 2025/1223.
\url{https://eprint.iacr.org/2025/1223}.

\bibitem{pcfsparse2025}
Sebastian Hasler and Pascal Reisert.
\emph{Pseudorandom Correlation Functions for Multiparty Beaver Triples from Sparse LPN}.
Preprint, Cryptology ePrint Archive, Paper 2025/2002, 2025.
\url{https://eprint.iacr.org/2025/2002}.

\bibitem{pcgf22025}
Peihan Miao, Alice Murphy, Akshayaram Srinivasan, and Max Tromanhauser.
\emph{Pseudorandom Correlation Generators for Multiparty Beaver Triples over $\mathbb{F}_2$}.
Cryptology ePrint Archive, Paper 2025/1182, 2025.
\url{https://eprint.iacr.org/2025/1182}.

\bibitem{batchauth2026}
Elette Boyle, Niv Gilboa, Matan Hamilis, Yuval Ishai, and Ariel Nof.
\emph{Fast PCGs for Batch-Authenticated Multiplication Triples}.
CRYPTO 2026, pp.~495--526.
\url{https://link.springer.com/chapter/10.1007/978-3-032-35418-1_16}.

\bibitem{mixedmode2025}
Wenhao Zhang, Xiao Wang, and Chenkai Weng.
\emph{Efficient Mixed-Mode Oblivious RAMs}.
Cryptology ePrint Archive, Paper 2025/1057, 2025.
\url{https://eprint.iacr.org/2025/1057}.

\bibitem{ods2014}
Xiao Shaun Wang et al.
\emph{Oblivious Data Structures}.
ACM CCS 2014; full version, Cryptology ePrint Archive, Paper 2014/185.
\url{https://eprint.iacr.org/2014/185}.

\bibitem{cuckoo2004}
Rasmus Pagh and Flemming Friche Rodler.
\emph{Cuckoo Hashing}.
Journal of Algorithms, 51(2):122--144, 2004.
\url{https://www.rasmuspagh.net/papers/cuckoo-jour.pdf}.

\bibitem{mimc2016}
Martin Albrecht, Lorenzo Grassi, Christian Rechberger, Arnab Roy, and Tyge Tiessen.
\emph{MiMC: Efficient Encryption and Cryptographic Hashing with Minimal Multiplicative Complexity}.
ASIACRYPT 2016; Cryptology ePrint Archive, Paper 2016/492.
\url{https://eprint.iacr.org/2016/492}.

\bibitem{poseidon2021}
Lorenzo Grassi, Dmitry Khovratovich, Christian Rechberger, Arnab Roy, and Markus Schofnegger.
\emph{Poseidon: A New Hash Function for Zero-Knowledge Proof Systems}.
USENIX Security, 2021.
\url{https://www.usenix.org/conference/usenixsecurity21/presentation/grassi}.

\bibitem{sorting2022}
Gilad Asharov et al.
\emph{Efficient Secure Three-Party Sorting with Applications to Data Analysis and Heavy Hitters}.
ACM CCS 2022, pp.~125--138.
\url{https://eprint.iacr.org/2022/1595}.

\bibitem{joins2024}
Gilad Asharov, Koki Hamada, Dai Ikarashi, Ryo Kikuchi, Ariel Nof, Benny Pinkas, and Junichi Tomida.
\emph{Secure Statistical Analysis on Multiple Datasets: Join and Group-By}.
Cryptology ePrint Archive, Paper 2024/141, 2024;
major revision of the ACM CCS 2023 paper, with a revised author list.
\url{https://eprint.iacr.org/2024/141}.

\bibitem{compaction2020}
Gilad Asharov, Ilan Komargodski, Wei-Kai Lin, Enoch Peserico, and Elaine Shi.
\emph{Oblivious Parallel Tight Compaction}.
ITC 2020; Cryptology ePrint Archive, Paper 2020/125.
\url{https://eprint.iacr.org/2020/125}.

\end{thebibliography}

\end{document}
